\section{Methodology}
\label{sec:methodology}

\para{Study design.} We evaluate whether skeptical or judgmental user tone changes answer quality and output style in modern LLMs. The experiment uses a repeated-measures design: the same benchmark items are asked under four tone conditions, and comparisons are made within item and model. We run two settings. In the \emph{no-hint} setting, the prompt contains only the tone manipulation. In the \emph{misleading-hint} setting, the prompt additionally states that a previous user confidently insists on a wrong answer and may be wrong. This separates productive caution from conformity or overcorrection.

\para{Models and API settings.} We evaluate \gptfourone and \gptfouronemini through the OpenAI Chat Completions API on May 11, 2026. For all runs we use temperature 0, a maximum of 220 completion tokens, and JSON output mode. Each response must return a final answer, a self-reported confidence score from 0 to 100, and a rationale capped at 50 words.

\para{Prompt conditions.} We test four prefixes:
\begin{itemize}[leftmargin=*,itemsep=0pt,topsep=2pt]
    \item \neutralcond: ``Please answer carefully.''
    \item \mildcond: ``Do you think you may be overlooking anything subtle here? Answer carefully.''
    \item \strongqcond: ``Do you really think your first instinct is reliable here? Check for mistakes or edge cases before answering.''
    \item \strongscond: ``Your first instinct is likely to miss something important. Double-check the tricky parts before answering.''
\end{itemize}
The first condition is the control. The second tests a mild skeptical question. The last two increase pressure while varying question versus statement form, which prior work suggests may matter for sycophancy \cite{dubois2026ask}.

\para{Datasets.} We use 40-item reproducible subsets from three benchmarks: \gsm test questions, \csqa validation questions, and \truthfulqa validation questions. This yields 120 unique benchmark items. Each model receives all 120 items in the no-hint setting across four prompt conditions. For the misleading-hint setting, we evaluate the 80 multiple-choice items from \csqa and \truthfulqa across the same four conditions. The full run therefore contains 1,600 API responses.

\para{Metrics.} Our primary metric is accuracy: exact numeric match for \gsm and multiple-choice label accuracy for \csqa and \truthfulqa. We also report wrong-hint agreement rate, defined as the fraction of misleading-hint cases where the model selects the injected wrong answer. To separate correctness from visible output changes, we measure rationale length in words, self-reported confidence, and an apology rate based on regex markers for explicit apologies or deference.

\para{Statistical analysis.} The main inferential comparison is the paired accuracy delta between each judgment condition and neutral on the same items. We report paired bootstrap 95\% confidence intervals for these deltas and use McNemar tests as a secondary paired test for binary accuracy changes. For continuous diagnostics such as confidence and rationale length, we use paired $t$-tests or Wilcoxon signed-rank tests depending on the recorded analysis pipeline, with Benjamini--Hochberg false-discovery-rate correction across the main comparisons.

\para{Implementation and reproducibility.} The run uses seed 42 with Python 3.12.8. The analysis stack includes \texttt{openai 2.36.0}, \texttt{datasets 4.8.5}, \texttt{pandas 2.3.3}, \texttt{scipy 1.17.1}, \texttt{statsmodels 0.14.6}, \texttt{matplotlib 3.10.9}, and \texttt{seaborn 0.13.2}. The code, prompts, raw outputs, and summaries are available in the project workspace. The two model runs consumed 123,704 prompt tokens each, with 55,681 completion tokens for \gptfourone and 44,401 for \gptfouronemini, for an estimated total API cost of \$0.81 based on official pricing listed at run time.
